% ===== 3. Тема 1 =====
\begin{frame}{Тема 1. Цифровий тренер для автоспорту}
  \begin{columns}[T,onlytextwidth]
    \begin{column}{0.48\textwidth}
      \textbf{\color{accent}Призначення}\\
      Вебзастосунок: GPS-запис з камери чи телефона $\to$ розбір заїзду

      \medskip
      \textbf{\color{accent}Для кого}\\
      Новачки й аматори без тренера: прокатний картинг, трек-дні

      \medskip
      \textbf{\color{accent}Проблема}\\
      Немає об'єктивного зворотного зв'язку, тренер і телеметрія дорогі
    \end{column}
    \begin{column}{0.48\textwidth}
      \textbf{\color{accent}Основні можливості}
      \begin{itemize}
        \item Автоматичний поділ заїзду на кола
        \item Траєкторія на карті траси з межами
        \item Повороти, де втрачається найбільше часу
        \item 2–3 поради простою мовою з відео
      \end{itemize}
    \end{column}
  \end{columns}

  \note{[Доповідач 1] Перша тема — цифровий тренер для автоспорту. Уявіть новачка, який регулярно їздить на прокатний картинг. Він відчуває, що десь втрачає час, але не знає, де саме і що змінити. Живий тренер дорогий, а професійна телеметрія розрахована на досвідчених пілотів. Наш застосунок бере запис з екшн-камери GoPro чи трек із телефона, сам ділить заїзд на кола, накладає траєкторію на реальні межі траси і знаходить повороти, де пілот втрачає найбільше часу. На виході — не графіки, а дві-три поради простою мовою з цифрами і фрагментом відео, плюс план на наступний заїзд.}
\end{frame}

% ===== 4. Тема 2 =====
\begin{frame}{Тема 2. Розумний кошик для товарів}
  \begin{columns}[T,onlytextwidth]
    \begin{column}{0.48\textwidth}
      \textbf{\color{accent}Призначення}\\
      Окремий застосунок сам складає кошик у Сільпо через MCP

      \medskip
      \textbf{\color{accent}Для кого}\\
      Зайняті люди 22–40 років у великих містах, клієнти Сільпо

      \medskip
      \textbf{\color{accent}Проблема}\\
      Планувати їжу й покупки щодня — рутина, продукти викидають
    \end{column}
    \begin{column}{0.48\textwidth}
      \textbf{\color{accent}Основні можливості}
      \begin{itemize}
        \item Оцінка того, що є вдома, за історією покупок
        \item План харчування на 2–3 дні
        \item Автоматичний кошик із поясненнями
        \item Підтвердження, доставка чи самовивіз
      \end{itemize}
    \end{column}
  \end{columns}

  \note{[Доповідач 2] Друга тема — розумний кошик для товарів. Ідея така: застосунок сам складає замовлення в Сільпо на найближчі дні. Він знає, що людина зазвичай купує, за історією покупок оцінює, що вже закінчилося вдома, і пропонує простий план харчування на два-три дні. Увечері людина отримує готовий кошик, де біля кожного товару написана причина. Їй лишається переглянути його і обрати доставку чи самовивіз. Працює це через офіційний MCP-сервер, який Сільпо відкрило в серпні 2026 року. Аудиторія — зайняті люди 22–40 років у великих містах, які вже купують у Сільпо.}
\end{frame}
